\documentclass[12pt,a4paper]{report}

\usepackage[margin=1in]{geometry}
\usepackage{times}
\usepackage{setspace}
\usepackage{graphicx}
\usepackage{booktabs}
\usepackage{hyperref}
\usepackage{amsmath}
\usepackage{listings}
\usepackage{xcolor}
\usepackage{fancyhdr}
\usepackage{titlesec}
\usepackage{tocloft}
\usepackage{caption}

\onehalfspacing

\pagestyle{fancy}
\fancyhf{}
\fancyhead[L]{\leftmark}
\fancyhead[R]{\thepage}
\renewcommand{\headrulewidth}{0.4pt}

\titleformat{\chapter}[display]
  {\normalfont\Large\bfseries\centering}
  {\chaptertitlename\ \thechapter}{12pt}{\LARGE}

\lstset{
  basicstyle=\ttfamily\small,
  breaklines=true,
  frame=single,
  language=Python,
  keywordstyle=\color{blue},
  commentstyle=\color{gray},
}

\begin{document}

% ========== TITLE PAGE ==========
\begin{titlepage}
\begin{center}
\vspace*{0.5cm}
\textbf{\Large INSTITUTE OF AERONAUTICAL ENGINEERING (IARE)}\\[0.3cm]
\textit{(Autonomous)}\\
Dundigal, Hyderabad -- 500 043\\[1cm]

\textbf{\large Department of Computer Science and Engineering}\\[2cm]

\textbf{\LARGE PROJECT REPORT}\\[0.5cm]
\textit{on}\\[0.5cm]

\textbf{\huge AI Agent Briefer}\\[0.3cm]
\textbf{\Large A Serverless, Multi-Agent System for\\Autonomous Daily Intelligence Synthesis}\\[2cm]

\textbf{Submitted by:}\\[0.3cm]
\begin{tabular}{ll}
Ramadugu Dhanush & ramadugudhanush8@gmail.com \\
\end{tabular}\\[1.5cm]

\textbf{Under the Guidance of:}\\[0.3cm]
\textbf{Mr. Gautham Sharma}\\
Assistant Professor\\
Department of Computer Science and Engineering\\[1.5cm]

\textbf{Academic Year 2025--2026}
\end{center}
\end{titlepage}

% ========== CERTIFICATE ==========
\newpage
\thispagestyle{empty}
\begin{center}
\textbf{\Large CERTIFICATE}\\[1.5cm]
\end{center}

This is to certify that the project report entitled \textbf{``AI Agent Briefer: A Serverless, Multi-Agent System for Autonomous Daily Intelligence Synthesis''} is a bonafide record of work carried out by \textbf{Ramadugu Dhanush}, in partial fulfillment of the requirements for the award of the degree of \textbf{Bachelor of Technology} in \textbf{Computer Science and Engineering} from the \textbf{Institute of Aeronautical Engineering (IARE)}, Hyderabad, during the academic year \textbf{2025--2026}.\\[3cm]

\noindent
\begin{minipage}{0.45\textwidth}
\centering
\rule{5cm}{0.4pt}\\
\textbf{Mr. Gautham Sharma}\\
Project Guide\\
Dept. of CSE
\end{minipage}
\hfill
\begin{minipage}{0.45\textwidth}
\centering
\rule{5cm}{0.4pt}\\
\textbf{Head of Department}\\
Dept. of CSE
\end{minipage}

% ========== ACKNOWLEDGMENT ==========
\newpage
\thispagestyle{empty}
\begin{center}
\textbf{\Large ACKNOWLEDGMENT}\\[1cm]
\end{center}

I would like to express my sincere gratitude to my project guide, \textbf{Mr. Gautham Sharma}, Assistant Professor, Department of Computer Science and Engineering, for his invaluable guidance, encouragement, and support throughout the duration of this project.

I extend my thanks to the \textbf{Head of the Department} and all the \textbf{faculty members} of the Department of Computer Science and Engineering, IARE, for providing the necessary infrastructure and resources.

I also thank my \textbf{family and friends} for their constant support and motivation during this project.\\[2cm]

\hfill \textbf{Ramadugu Dhanush}

% ========== ABSTRACT ==========
\newpage
\thispagestyle{empty}
\begin{center}
\textbf{\Large ABSTRACT}\\[1cm]
\end{center}

In the modern era, the rapid dissemination of information has led to an overwhelming volume of digital news, making efficient curation a significant challenge. Concurrently, technical interview preparation requires synthesizing vast amounts of decentralized corporate data, which is time-consuming and often yields outdated information.

This project presents ``AI Agent Briefer,'' an autonomous, serverless, multi-agent artificial intelligence system designed to deliver high-precision daily technology briefings and in-depth company intelligence. Built on the CrewAI orchestration framework, the system leverages a robust multi-model fallback chain---integrating Large Language Models (LLMs) such as Ollama, Groq, Gemini, OpenRouter, Mistral, and Together AI---to ensure high uptime with zero operational costs.

A critical innovation of this project is its rigorous programmatic output validation pipeline, which systematically identifies and mitigates LLM hallucinations, lazy tool execution, and stale data generation through nine distinct deterministic checks. The curated intelligence is delivered via Telegram acting as a multi-modal interface, providing both rich HTML text and synthesized audio via ElevenLabs and Google TTS.

The system is deployed on a fully serverless infrastructure combining Vercel Functions, GitHub Actions, and the Telegram Bot API, achieving zero hosting costs and zero infrastructure management.

\textbf{Keywords:} Multi-Agent Systems, CrewAI, Large Language Models, Serverless Architecture, Output Validation, Generative AI, Text-to-Speech.

% ========== TABLE OF CONTENTS ==========
\tableofcontents
\listoftables

% ========== CHAPTER 1: INTRODUCTION ==========
\chapter{Introduction}

\section{Background}
The technology industry is characterized by rapid and relentless evolution. Professionals, job seekers, and technology enthusiasts must continuously monitor advancements across Artificial Intelligence, cybersecurity, cloud computing, startups, and global markets. General-purpose news aggregators such as Google News or Apple News lack the domain specificity required to filter irrelevant content, while manual research for technical interview preparation is time-consuming and often yields outdated information.

The advent of autonomous AI agents---software entities capable of reasoning, planning, using external tools, and synthesizing information---presents a profound opportunity to automate this knowledge curation process. Frameworks such as CrewAI and AutoGen enable developers to define specialized agent personas that can independently search the web, analyze results, and produce structured reports without continuous human supervision.

\section{Problem Statement}
Existing AI-driven summarization tools are hindered by three critical limitations:
\begin{enumerate}
    \item \textbf{High API Costs:} Reliance on proprietary models such as OpenAI's GPT-4 prevents continuous daily execution without significant financial investment. A single daily briefing with web search can cost \$1--2 per run, rendering automated daily operations financially unsustainable.
    \item \textbf{LLM Hallucinations and Functional Laziness:} Generative models frequently fabricate dates, summarize outdated articles as ``breaking news,'' or fail to execute their assigned web search tools---instead generating responses like ``I could not find this information'' without actually attempting the search.
    \item \textbf{Limited Accessibility:} Text-heavy newsletters and reports are difficult to consume in varying environments (commuting, exercising, multitasking), highlighting the need for multi-modal delivery systems that include audio synthesis.
\end{enumerate}

\section{Project Objectives}
The primary objectives of this project are:
\begin{enumerate}
    \item To develop distinct AI personas capable of specialized tasks: a \textit{News Scout} agent for daily technology briefings and an \textit{Interview Strategist} agent for comprehensive company intelligence.
    \item To engineer a zero-cost inference engine by implementing a dynamic, rate-limit-aware fallback mechanism across seven free-tier LLM providers.
    \item To design deterministic Python guardrails that programmatically reject subpar generative outputs through nine distinct validation checks.
    \item To deploy the system as a fully serverless infrastructure utilizing GitHub Actions (cron scheduling), Vercel Functions (webhook handling), and the Telegram Bot API (user interaction) for seamless multi-modal delivery via text and audio.
\end{enumerate}

\section{Scope and Contributions}
This work makes the following key technical contributions:
\begin{enumerate}
    \item A multi-model fallback architecture traversing up to seven LLM providers with rate-limit-aware exponential backoff, achieving theoretically unbounded uptime at zero cost.
    \item A programmatic output validation pipeline with nine deterministic checks covering stale dates, headline deduplication, topic diversity, lazy agent detection, tool-call leakage, unfilled placeholders, geographic balance, meta-message contamination, and minimum content thresholds.
    \item A three-tier serverless deployment architecture combining Vercel Functions, GitHub Actions, and the Telegram Bot API for zero-infrastructure operation.
    \item A dual text-to-speech integration (ElevenLabs + Google gTTS) with Markdown-to-speech pre-processing for multi-modal delivery.
    \item Constraint-driven prompt engineering techniques that programmatically combat LLM laziness through mandated search counts, coverage ratios, and output schemas.
\end{enumerate}

\section{Organization of the Report}
The remainder of this report is organized as follows:
\begin{itemize}
    \item \textbf{Chapter 2} presents a literature review covering multi-agent systems, LLM hallucinations, serverless computing, TTS synthesis, and cost-optimized inference.
    \item \textbf{Chapter 3} describes the system architecture, deployment infrastructure, technology stack, and security design.
    \item \textbf{Chapter 4} details the implementation of the LLM fallback chain, agent personas, webhook handler, audio pipeline, message sanitization, rate-limit handling, and output post-processing.
    \item \textbf{Chapter 5} presents the validation and testing framework with nine guardrail categories.
    \item \textbf{Chapter 6} evaluates cost efficiency, execution latency, validation effectiveness, model comparison, and end-to-end reliability.
    \item \textbf{Chapter 7} concludes the report with lessons learned and future scope.
\end{itemize}

% ========== CHAPTER 2: LITERATURE REVIEW ==========
\chapter{Literature Review}

\section{Autonomous Multi-Agent Systems}
Multi-agent systems (MAS) involve multiple interacting intelligent agents that solve problems beyond the capability of a single agent. Frameworks like CrewAI and AutoGen have popularized the use of LLMs as the cognitive engine for these agents, allowing them to collaborate, delegate tasks, and utilize external tools (e.g., search APIs) to gather real-time data autonomously.

The CrewAI framework provides a high-level abstraction for defining agent roles, goals, backstories, and tool bindings. It manages the conversation loop between the LLM ``brain'' and external tools, handling the ReAct (Reasoning + Acting) paradigm internally. Unlike AutoGen's multi-turn dialogue approach, CrewAI emphasizes task-oriented sequential and hierarchical execution patterns, which are better suited for deterministic, pipeline-oriented workloads like news curation.

\section{LLM Hallucinations and Output Validation}
A well-documented limitation of LLMs is ``hallucination''---the generation of fluent but factually incorrect information. Ji et al. provide a comprehensive taxonomy categorizing hallucinations into intrinsic (contradicting the source) and extrinsic (unverifiable against any source) types. In automated news curation, this manifests as stale news being presented as recent.

A particularly pernicious form in agentic workflows is \textit{functional laziness}---where the LLM fabricates a response rather than invoking its assigned tool to retrieve actual data. Recent studies suggest that combining stochastic LLM generation with deterministic programmatic validation (e.g., regex-based date validation) significantly improves the reliability of enterprise AI systems.

\section{Serverless Computing}
Serverless computing allows developers to execute code without managing underlying infrastructure. The Function-as-a-Service (FaaS) paradigm employed by Vercel executes individual HTTP handler functions on demand. GitHub Actions extends this with workflow-based scheduled cron triggers, manual dispatches, and event triggers---within a free tier of 2,000 minutes per month. This combination enables complex multi-stage pipelines without provisioning persistent compute resources.

\section{Multi-Modal Content Delivery}
Text-to-speech synthesis has evolved from concatenative approaches to neural TTS models producing near-human voice quality. ElevenLabs provides state-of-the-art neural TTS with configurable voice identities. Google's gTTS library provides a lightweight, free-tier alternative. This project targets condensed audio summaries optimized for mobile consumption via Telegram's native audio player, requiring Markdown stripping and character-count truncation.

\section{Cost-Optimized AI Inference}
Commercial API pricing (e.g., OpenAI GPT-4o at \$5.00 per million input tokens) renders daily automated workloads financially unsustainable. Free-tier offerings from Groq (unlimited), Google AI Studio (15 req/min), and OpenRouter (limited free quotas) create opportunities for cost-aware inference routing. LiteLLM provides a unified API abstraction across 100+ providers, enabling seamless provider switching through model string prefixing.

% ========== CHAPTER 3: SYSTEM ARCHITECTURE ==========
\chapter{System Architecture and Design}

\section{Overview}
The system is designed with a decoupled architecture, separating orchestration, inference, validation, and delivery pipelines to ensure module independence and high cohesion.

\section{Component Architecture}
The system comprises five major components:
\begin{enumerate}
    \item \textbf{User Interface (Telegram Bot):} Acts as the primary interface. Users receive automated daily briefs or request on-demand interview prep via \texttt{/brief} or \texttt{/interview <company>}.
    \item \textbf{Orchestrator (\texttt{main.py}):} The central coordination module. It initializes tasks, sets up the CrewAI environment, manages the LLM fallback loop, invokes validation, and dispatches results.
    \item \textbf{Agent Factory (\texttt{agents.py}):} Defines agent personas, equips them with tools (Serper Search API), manages LLM instantiation, and provides TTS and Telegram delivery functions.
    \item \textbf{Validation Engine:} A suite of deterministic functions performing nine distinct quality checks on raw LLM output before permitting dispatch.
    \item \textbf{Delivery and TTS Subsystem:} Formats validated text into Telegram-compatible HTML, generates MP3 audio files, and handles message chunking for large outputs.
\end{enumerate}

\section{Data Flow}
The system follows a four-stage pipeline:
\begin{enumerate}
    \item \textbf{Trigger:} A GitHub Actions cron job fires at 6:00 AM IST daily, or a Telegram webhook command triggers execution on demand.
    \item \textbf{Inference:} The Orchestrator requests an LLM from the Agent Factory. The AI agent receives its structured prompt and uses the Serper tool to fetch live web data.
    \item \textbf{Validation:} The Orchestrator parses the raw markdown response through the validation engine. If the output fails any check, the system throws an exception and retries with the next LLM in the fallback chain.
    \item \textbf{Delivery:} The validated text is sanitized into Telegram HTML, sent as a message, and passed to the TTS subsystem for audio synthesis and dispatch.
\end{enumerate}

\section{Serverless Deployment Architecture}
The production deployment uses a three-tier serverless pipeline:

\textbf{Tier 1 --- Webhook Ingestion (Vercel Functions):} The Telegram Bot API forwards all messages to a Vercel-hosted serverless function at the \texttt{/webhook} endpoint. The handler (\texttt{api/webhook.py}) parses incoming JSON, deduplicates webhooks using an in-memory timestamp cache (60-second expiry), routes commands, and triggers GitHub Actions workflows via the \texttt{workflow\_dispatch} API.

\textbf{Tier 2 --- Compute Orchestration (GitHub Actions):} The Vercel function triggers the \texttt{briefer.yml} workflow, which provisions an Ubuntu runner, installs Python 3.11 and dependencies, verifies critical imports, and executes the Python entry point with all API keys injected from GitHub Secrets.

\textbf{Tier 3 --- Scheduled Execution (Cron):} GitHub Actions provides native cron scheduling via \texttt{cron: '30 0 * * *'} (00:30 UTC / ~6:00 AM IST), executing daily briefings without user intervention.

\section{Technology Stack}

\begin{table}[htbp]
\caption{Technology Stack Summary}
\centering
\begin{tabular}{@{}llp{5cm}@{}}
\toprule
\textbf{Layer} & \textbf{Technology} & \textbf{Role} \\
\midrule
Orchestration & CrewAI & Agent definition, task execution, tool binding \\
LLM Routing & LiteLLM & Unified API across 7 providers \\
Web Search & Serper API & Real-time web and news search \\
TTS (Primary) & ElevenLabs & Neural voice synthesis \\
TTS (Fallback) & Google gTTS & Free-tier voice synthesis \\
Messaging & Telegram Bot API & User interface and delivery \\
Hosting & Vercel Functions & Serverless webhook handler \\
CI/CD & GitHub Actions & Cron scheduling, workflow dispatch \\
HTTP & \texttt{requests} & API calls, Telegram dispatch \\
Config & \texttt{python-dotenv} & Environment variable management \\
SSL & \texttt{certifi} & Certificate verification \\
\bottomrule
\end{tabular}
\end{table}

\section{Security Architecture}
\begin{itemize}
    \item \textbf{Secret Isolation:} All API keys are stored in \texttt{.env} (local), GitHub Secrets (CI/CD), or Vercel Environment Variables (webhook). The \texttt{.gitignore} excludes \texttt{.env} from version control.
    \item \textbf{Least Privilege:} The GitHub PAT is scoped to only \texttt{repo} and \texttt{workflow} permissions with 90-day expiration.
    \item \textbf{Chat Restriction:} \texttt{TELEGRAM\_CHAT\_ID} restricts automated delivery; on-demand commands dynamically route to the sender's chat.
    \item \textbf{Input Sanitization:} HTML entities (\texttt{\&}, \texttt{<}, \texttt{>}) are escaped before formatting to prevent injection via LLM output.
\end{itemize}

% ========== CHAPTER 4: IMPLEMENTATION ==========
\chapter{Implementation Details}

\section{The LLM Fallback Chain}
The \texttt{get\_llm()} factory function dynamically routes inference across providers in a configurable priority order:
\begin{enumerate}
    \item \textbf{Ollama (Local):} Zero latency, zero API limits.
    \item \textbf{Groq (Llama 3.3 70B):} Extreme inference speed, generous free tier.
    \item \textbf{Google Gemini (2.0 Flash):} 2-million-token context window for search-heavy tasks.
    \item \textbf{Together AI, HuggingFace, OpenRouter, Mistral:} Progressive safety nets.
\end{enumerate}

The function accepts \texttt{skip\_models} (providers to exclude during fallback) and \texttt{prefer\_model} (to force a specific provider, e.g., Gemini for interview tasks). Each provider is instantiated via LiteLLM's model string convention with provider-specific API keys, 120-second timeouts, temperature 0.7, and 2 max retries.

\section{Agent Personas and Prompt Engineering}
Two specialized agents are defined via factory functions:

\textbf{The News Scout} (\texttt{create\_news\_scout\_agent()}): Role ``Tech News Analyst'' with a dynamically injected date anchor. The task prompt mandates exactly 5 topic-specific searches (AI/ML, cybersecurity, startup funding, India tech, global tech), each including the current date string. Configuration: \texttt{max\_iter=12}, \texttt{max\_execution\_time=600s}, \texttt{allow\_delegation=False}.

\textbf{The Interview Strategist} (\texttt{create\_company\_researcher\_agent()}): Executes a 4-pronged search methodology (Overview, News, Culture, Challenges). The prompt commands generation of exactly 5 ``Killer Questions'' referencing specific facts from search results. Configuration: \texttt{max\_iter=12}, \texttt{max\_execution\_time=480s}.

\section{Webhook Event-Driven Architecture}
The \texttt{api/webhook.py} module implements the serverless Telegram command handler on Vercel. It supports three commands: \texttt{/brief} (triggers daily briefing workflow), \texttt{/interview <company>} (triggers interview prep), and \texttt{/help} (returns usage instructions).

A critical reliability feature is the duplicate webhook deduplication cache---an in-memory dictionary mapping \texttt{update\_id} values to timestamps with 60-second expiry, preventing double-execution when Telegram resends slow-responding webhooks. GitHub Actions workflows are triggered via the \texttt{workflow\_dispatch} REST API with Bearer token authentication.

\section{Audio Synthesis Pipeline}
The TTS subsystem converts validated briefings into MP3 audio:
\begin{enumerate}
    \item \textbf{Text Pre-processing:} \texttt{\_extract\_audio\_summary()} splits on the 📰 delimiter to extract headline items; \texttt{\_strip\_markdown\_for\_speech()} removes all HTML tags, Markdown syntax, bullet markers, and formatting artifacts to produce speech-friendly text.
    \item \textbf{ElevenLabs (Primary):} Uses the \texttt{eleven\_turbo\_v2\_5} model with configurable voice selection via \texttt{\_pick\_voice\_id()}. Default: ``Rachel'' voice.
    \item \textbf{Google gTTS (Fallback):} Configured with Australian English accent (\texttt{tld="com.au"}) for natural-sounding output.
    \item \textbf{Dispatch:} The MP3 file is sent via Telegram's \texttt{sendAudio} API endpoint.
\end{enumerate}

\section{Telegram Message Sanitization}
The \texttt{format\_for\_telegram()} function operates as a 7-stage Markdown-to-Telegram-HTML transpiler:
\begin{enumerate}
    \item HTML entity escaping (\texttt{\&}, \texttt{<}, \texttt{>})
    \item Horizontal rule removal
    \item Table transpilation (rows → bulleted key-value pairs)
    \item Header conversion (\texttt{\#\#} → \texttt{<b>})
    \item Inline formatting (bold, italic, code)
    \item Whitespace normalization
    \item Message chunking along paragraph boundaries (4000-char limit with 96-char buffer)
\end{enumerate}

\section{Rate Limit Handling and Exponential Backoff}
The system employs two-tiered resilience:

\textbf{Intra-model retry:} \texttt{run\_crew\_with\_rate\_limit\_retry()} allows up to 3 retries per model. Wait time follows:
\begin{equation}
W = W_{\text{base}} + (\text{attempt} \times 45)
\end{equation}
where $W_{\text{base}}$ is parsed from the 429 error payload (default: 90s), producing waits of ~90s, 135s, 180s.

\textbf{Inter-model fallback:} The orchestrator applies progressive delays between model switches:
\begin{equation}
D = (\text{attempt} - 1) \times 30
\end{equation}
producing 0s, 30s, 60s, 90s, 120s across five attempts---yielding up to 15 total inference attempts.

\section{Interview Output Post-Processing}
The \texttt{clean\_interview\_preamble()} function strips LLM ``thinking aloud'' preambles by locating the first section header pattern and removing preceding content exceeding 50 characters. It also iteratively removes trailing meta-commentary phrases (e.g., ``the final answer'', ``do not copy these instructions'').

% ========== CHAPTER 5: VALIDATION ==========
\chapter{Validation and Testing}

\section{Overview}
The most vulnerable point of autonomous agents is output reliability. This project implements nine distinct deterministic validation checks on raw LLM output before dispatch.

\section{News Validation Guardrails}

\subsection{Stale News Filtering}
Extracts date strings in both ISO (\texttt{YYYY-MM-DD}) and natural language formats (``March 18, 2026'', ``Mar 18th, 2026''), including ordinal suffixes. Converts to \texttt{datetime} objects and rejects stories exceeding \texttt{NEWS\_MAX\_AGE\_DAYS} (default: 3 days).

\subsection{Headline Deduplication}
Computes word-overlap similarity between all headline pairs. If $\frac{|A \cap B|}{\max(|A|, |B|)} > 0.60$ for any two headlines, validation fails.

\subsection{Topic Diversity Enforcement}
Maintains two keyword arrays: consumer keywords (12 terms: ``launch'', ``smartphone'', ``specs'', etc.) and diversity keywords (17 terms: ``ai'', ``cybersecurity'', ``funding'', etc.). If no diversity keywords appear, validation rejects the output.

\subsection{Geographic Coverage}
Uses dual keyword sets: 22 India-specific terms (cities, companies like ``infosys'', ``flipkart'', ``upi'') and 18 global terms (``google'', ``nvidia'', ``silicon valley''). Both must match for validation to pass.

\subsection{Minimum Story Count}
Requires at least 5 unique headlines (identified by the 📰 emoji prefix) after deduplication.

\section{Interview Validation Guardrails}

\subsection{Anti-Lazy Agent Detection}
If ``Not found in public sources'' appears 4+ times, the agent is assumed to have skipped tool execution. Validation fails with a diagnostic message.

\subsection{Tool-Call Leakage Prevention}
Detects 11 raw tool-call patterns (e.g., \texttt{search\_the\_internet\_with\_serper}, \texttt{action\_input}, \texttt{<search\_query>}). If 2+ patterns are found, the output is rejected as containing debug syntax.

\subsection{Unfilled Placeholder Detection}
Regex scans for bracket placeholders like \texttt{[year]}, \texttt{[Name]}, \texttt{[Framework]}. If 3+ are found, the output is classified as an unfilled template.

\section{Common Guardrails}

\subsection{Meta-Message Detection}
Checks for 6 meta-phrases (``let's search'', ``i need to'', ``call another tool'', etc.) in outputs shorter than 1000 characters, catching LLM ``thinking aloud'' responses.

\subsection{Minimum Content Length}
Both validators reject outputs under 200 characters. The orchestrator independently enforces a 100-character minimum on raw crew results.

% ========== CHAPTER 6: RESULTS ==========
\chapter{Results and Performance Evaluation}

\section{Cost Efficiency}
By routing through free-tier providers, operational cost is \textbf{\$0.00/month}. An equivalent GPT-4o system would cost an estimated \$40/month.

\begin{table}[htbp]
\caption{Monthly Cost Comparison}
\centering
\begin{tabular}{@{}lrr@{}}
\toprule
\textbf{Component} & \textbf{This System} & \textbf{GPT-4o Equiv.} \\
\midrule
LLM Inference (30 runs) & \$0.00 & \$30.00 \\
Web Search API & \$0.00 & \$5.00 \\
TTS Audio & \$0.00 & \$3.00 \\
Compute (CI/CD) & \$0.00 & \$2.00 \\
\midrule
\textbf{Total} & \textbf{\$0.00} & \textbf{\$40.00} \\
\bottomrule
\end{tabular}
\end{table}

\section{Execution Latency}
\begin{itemize}
    \item \textbf{Daily Briefing:} 25--45 seconds (bound by Serper API network latency).
    \item \textbf{Interview Prep:} 40--80 seconds (Gemini's 2M-token context processes dozens of pages without chunking).
    \item \textbf{Audio Synthesis:} ElevenLabs generates 2-minute MP3s in under 10 seconds; gTTS fallback adds ~5 seconds.
\end{itemize}

\section{Validation Effectiveness}

\begin{table}[htbp]
\caption{Observed Validation Failure Categories}
\centering
\begin{tabular}{@{}lp{5cm}l@{}}
\toprule
\textbf{Failure Type} & \textbf{Description} & \textbf{Frequency} \\
\midrule
Stale News & Stories older than 3 days & High \\
Lazy Agent & ``Not found'' responses & Medium \\
Tool-Call Leakage & Raw XML/JSON in output & Medium \\
Missing Dates & No publication dates & Medium \\
Deduplication & $>$60\% headline overlap & Low \\
Placeholders & Unfilled \texttt{[brackets]} & Low \\
Meta-Messages & ``I need to search'' & Low \\
Topic Monotonicity & Only product launches & Rare \\
Geographic Imbalance & Missing India/Global & Rare \\
\bottomrule
\end{tabular}
\end{table}

\section{Model Performance Comparison}

\begin{table}[htbp]
\caption{LLM Provider Performance Comparison}
\centering
\begin{tabular}{@{}llll@{}}
\toprule
\textbf{Provider} & \textbf{Speed} & \textbf{Context} & \textbf{Free Tier} \\
\midrule
Groq & Very Fast & 128K & Unlimited \\
Gemini Flash & Fast & 2M & 15 req/min \\
Together AI & Fast & 128K & \$25 credits \\
HuggingFace & Moderate & 8K & 1K req/day \\
OpenRouter & Fast & 128K & Limited \\
Mistral & Fast & 32K & Limited \\
\bottomrule
\end{tabular}
\end{table}

\section{End-to-End Reliability}
\begin{itemize}
    \item \textbf{Inference:} With seven providers, simultaneous rate-limit exhaustion is extremely unlikely. Stress testing with invalidated keys confirmed 100\% request fulfillment via fallback.
    \item \textbf{Audio:} ElevenLabs' 10K chars/month free tier covers ~5 briefings; gTTS fallback ensures uninterrupted audio delivery.
    \item \textbf{Chunking:} Interview outputs regularly exceed 4096 characters; paragraph-boundary splitting produces 2--4 coherent message chunks.
    \item \textbf{Error Notification:} After all models are exhausted, structured error messages with actionable troubleshooting are sent to the user via Telegram.
\end{itemize}

% ========== CHAPTER 7: CONCLUSION ==========
\chapter{Conclusion and Future Scope}

\section{Conclusion}
The ``AI Agent Briefer'' successfully demonstrates the viability of deploying highly autonomous knowledge-curation agents on strictly serverless, zero-cost infrastructure. By acknowledging the innate flaws of modern LLMs---hallucinations and functional laziness---and framing them within stringent programmatic control boundaries, the pipeline transforms stochastic textual generation into deterministic, enterprise-grade intelligence summaries accessible through Telegram.

\section{Lessons Learned}
\begin{enumerate}
    \item \textbf{Deterministic validation is non-negotiable for agentic AI.} The nine-check pipeline rejected approximately 15--25\% of raw LLM outputs during testing, each of which would have delivered incorrect content.
    \item \textbf{Fallback chains outperform single-model architectures.} Multi-model routing transforms provider failures into brief delays rather than system crashes.
    \item \textbf{Serverless matches periodic AI workloads.} A system running 60 seconds daily would leave a persistent server 99.93\% idle. GitHub Actions + Vercel perfectly matches this bursty pattern at zero cost.
    \item \textbf{Prompt engineering must be adversarial.} LLMs exhibit lazy behaviors with underspecified prompts. Effective agentic prompts require structural constraints (``Do exactly 5 searches'', ``Every story must include a date'').
\end{enumerate}

\section{Future Scope}
\begin{enumerate}
    \item \textbf{Vector Database Integration (RAG):} ChromaDB-based temporal memory to prevent cross-day news duplication.
    \item \textbf{Dynamic User Profiling:} Configurable topic preferences replacing hardcoded geographic parameters.
    \item \textbf{Hierarchical Agent Topologies:} Parallel execution layers to reduce latency by an estimated 50\%.
    \item \textbf{Cross-Platform Messaging:} Expansion to Slack, Discord, and WhatsApp Business API.
    \item \textbf{Multi-Language TTS:} Support for Hindi, Telugu, and other regional languages.
    \item \textbf{Real-Time Streaming:} SSE/WebSocket-based incremental output delivery.
    \item \textbf{Automated Quality Scoring:} LLM-as-judge secondary evaluation before dispatch.
\end{enumerate}

% ========== REFERENCES ==========
\begin{thebibliography}{99}

\bibitem{crewai} J.~M.~Moura, ``CrewAI: Framework for orchestrating role-playing, autonomous AI agents,'' GitHub Repository, 2024. Available: \url{https://github.com/joaomdmoura/crewAI}

\bibitem{litellm} BerriAI, ``LiteLLM: Call all LLM APIs using the OpenAI format,'' GitHub Repository, 2024. Available: \url{https://github.com/BerriAI/litellm}

\bibitem{llama2} H.~Touvron et al., ``Llama 2: Open foundation and fine-tuned chat models,'' \textit{arXiv:2307.09288}, 2023.

\bibitem{ji2023} Z.~Ji et al., ``Survey of hallucination in natural language generation,'' \textit{ACM Computing Surveys}, vol.~55, no.~12, 2023.

\bibitem{githubactions} GitHub, ``GitHub Actions documentation,'' 2024. Available: \url{https://docs.github.com/en/actions}

\bibitem{elevenlabs} ElevenLabs, ``ElevenLabs API documentation,'' 2024. Available: \url{https://elevenlabs.io/docs}

\bibitem{serper} Serper, ``Serper Search API documentation,'' 2024. Available: \url{https://serper.dev/}

\bibitem{vercel} Vercel, ``Vercel serverless functions documentation,'' 2024. Available: \url{https://vercel.com/docs/functions}

\bibitem{telegram} Telegram, ``Telegram Bot API documentation,'' 2024. Available: \url{https://core.telegram.org/bots/api}

\bibitem{gtts} P.~N.~Dekimpe, ``gTTS: Google Text-to-Speech,'' GitHub Repository, 2024. Available: \url{https://github.com/pndurette/gTTS}

\bibitem{autogen} Q.~Wu et al., ``AutoGen: Enabling next-gen LLM applications via multi-agent conversation,'' \textit{arXiv:2308.08155}, 2023.

\bibitem{react} S.~Yao et al., ``ReAct: Synergizing reasoning and acting in language models,'' \textit{arXiv:2210.03629}, 2023.

\end{thebibliography}

% ========== APPENDIX ==========
\appendix
\chapter{Source Code Repository}
The complete source code for this project is available at:\\
\url{https://github.com/RDK7159357/ai-agents-using-crewai}

\noindent Key files:
\begin{itemize}
    \item \texttt{main.py} --- Orchestrator, fallback logic, validation, Telegram delivery (766 lines)
    \item \texttt{agents.py} --- LLM factory, agent personas, TTS pipeline, messaging (442 lines)
    \item \texttt{api/webhook.py} --- Vercel serverless webhook handler (157 lines)
    \item \texttt{.github/workflows/briefer.yml} --- GitHub Actions workflow (85 lines)
    \item \texttt{vercel.json} --- Vercel deployment configuration
    \item \texttt{requirements.txt} --- Python dependencies (12 packages)
\end{itemize}

\chapter{Environment Configuration}
The system requires the following environment variables (all free-tier):
\begin{itemize}
    \item \textbf{LLM Providers (at least one):} \texttt{GROQ\_API\_KEY}, \texttt{GOOGLE\_API\_KEY}, \texttt{TOGETHER\_API\_KEY}, \texttt{HUGGINGFACE\_API\_KEY}, \texttt{OPENROUTER\_API\_KEY}, \texttt{MISTRAL\_API\_KEY}
    \item \textbf{Web Search:} \texttt{SERPER\_API\_KEY} (2,500 free searches/month)
    \item \textbf{Messaging:} \texttt{TELEGRAM\_BOT\_TOKEN}, \texttt{TELEGRAM\_CHAT\_ID}
    \item \textbf{Audio (Optional):} \texttt{ELEVENLABS\_API\_KEY} (10K chars/month free)
    \item \textbf{Configuration:} \texttt{PREFERRED\_MODEL}, \texttt{NEWS\_MAX\_AGE\_DAYS}, \texttt{AUDIO\_USE\_SUMMARY}
\end{itemize}

\end{document}
